\begin{abstract}
 DNA-encoded chemical libraries (DELs) enable the screening of millions to billions of DNA-tagged small molecules in a single experiment, with binders identified after PCR amplification and high-throughput DNA sequencing. Analysis of DEL selection data requires the integration of sequence processing, cheminformatics, and statistical analysis, yet software solutions supporting these steps are fragmented, limited in scope, or proprietary. This protocol describes DELT-Hit (DNA-Encoded Library Technology Hit Identification), an open-source command-line tool for end-to-end DEL data analysis. DELT-Hit converts raw FASTQ reads into analysis-ready chemical data through five connected modules: sequence demultiplexing with error-tolerant barcode handling; chemical structure reconstruction from building block libraries and reaction SMIRKS templates; molecular property and descriptor calculation; enrichment analysis and hit ranking using count-based and RNA-seq-derived methods; and quality-control and visualization tools for DEL data interpretation. The workflow supports both single-display and dual-display DEL architectures and produces standardized outputs for downstream analysis, including machine-learning applications. Basic command-line skills are sufficient for routine analysis, whereas extending library enumeration workflows requires cheminformatics expertise. Depending on library size and sequencing depth, the workflow typically requires one hour to several hours. DELT-Hit is available at \url{https://github.com/DELTechnology/delt-hit}.
\end{abstract}

\section*{Key points}

\begin{itemize}
  \item DELT-Hit processes large DEL datasets and supports libraries containing millions of compounds.
  \item The workflow supports both single-display and dual-display DEL architectures.
  \item Established tools such as Cutadapt and edgeR are integrated with DEL-specific sequence handling, enumeration, and quality control.
  \item The modular design supports diverse library formats, custom reaction templates, and user-defined building block tables.
  \item Standardized outputs can be used directly for downstream cheminformatics and machine-learning analyses.
\end{itemize}
